% =====================================================================
%  Präambel
% =====================================================================

% ---------- Sprache, Kodierung, einheitliche Schrift -----------------
\usepackage[T1]{fontenc}
\usepackage{babel}            % Sprache kommt aus der Klassenoption ngerman
\usepackage{lmodern}          % eine Schriftfamilie für das gesamte Dokument
\usepackage{microtype}
\usepackage[autostyle]{csquotes} % \enquote{...} setzt korrekte „Anführungszeichen“

% KOMA setzt Überschriften standardmäßig serifenlos -> auf Grundschrift umstellen,
% damit überall dieselbe Schriftart verwendet wird.
\setkomafont{disposition}{\normalfont\bfseries}
\setkomafont{caption}{\normalsize}
\setkomafont{captionlabel}{\normalsize\bfseries}

% ---------- Bilder und Tabellen --------------------------------------
\usepackage{graphicx}
\graphicspath{{bilder/}}
\usepackage{booktabs}         % saubere Tabellenlinien
\usepackage{tabularx}
\usepackage{array}
\usepackage{float}
\usepackage[section]{placeins}   % Bilder/Tabellen bleiben in ihrem Abschnitt

% ---------- Literatur: APA 7 ----------------------------------------
\usepackage[style=apa, backend=biber]{biblatex}
\addbibresource{literatur.bib}
\DefineBibliographyStrings{ngerman}{
  nodate      = {o.\,J.},% deutsch „ohne Jahr“
  retrieved   = {Abgerufen am},% deutsch „retrieved“
  from        = {von},% deutsch „at“
}

% ---------- Links und dynamische Verweise ----------------------------
\usepackage[hidelinks]{hyperref}
\usepackage{xurl}             % lange URLs umbrechen
\usepackage[ngerman, noabbrev]{cleveref} % \cref{...} -> „Abbildung 1“, „Tabelle 2“ …

% ---------- Hilfsbefehle ---------------------------------------------
% Quellenangabe für nicht selbst erstellte Bilder, z.\,B. in der Bildunterschrift:
%   \caption{Beschreibung. \bildquelle{schluessel}}
\newcommand{\bildquelle}[1]{(Quelle: \cite{#1})}

% Kennzeichnung von KI-/Automatisierungsnutzung direkt an der betroffenen Stelle:
%   \kihinweis{Tool}{Wofür genutzt}{bib-Schlüssel}
% Beispiel: \kihinweis{ChatGPT}{Gliederungsvorschlag, danach überarbeitet}{openai2026}
\newcommand{\kihinweis}[3]{%
  \footnote{KI-Nutzung: #1 -- #2 \parencite{#3}.}%
}

% Abschnitt „Beiträge“ am Ende jedes Kapitels:
%   \begin{beitraege}
%     Arbeitsschritt/Textteil & Verantwortlich \\
%   \end{beitraege}
\newenvironment{beitraege}{%
  \section{Beiträge}
  \noindent\tabularx{\linewidth}{@{}>{\raggedright\arraybackslash}X >{\raggedright\arraybackslash}p{0.3\linewidth}@{}}
  \toprule
  \textbf{Arbeitsschritt / Textanteil} & \textbf{Verantwortlich} \\
  \midrule
}{%
  \bottomrule
  \endtabularx
}
